% 2012 年江西卷文科第 7 题：主视图 3 个 1×1 方格、左视图 2 个、俯视图为六边形。
% 来源 PDF：普通高考/2012/2012江西文.pdf，第 1 页（印面 749）中栏。
% 六边形顶点 (1,2),(2,2),(3,1),(2,0),(1,0),(0,1)，面积 4；棱柱高 1，体积 4。
\begin{tikzpicture}[scale=0.95, semithick, font=\small]
  % 主视图
  \begin{scope}
    \draw (0,0) rectangle (3,1);
    \draw (1,0) -- (1,1); \draw (2,0) -- (2,1);
    \draw[<->] (0,-0.28) -- (1,-0.28); \node[below] at (0.5,-0.28) {1};
    \draw[<->] (1,-0.28) -- (2,-0.28); \node[below] at (1.5,-0.28) {1};
    \draw[<->] (2,-0.28) -- (3,-0.28); \node[below] at (2.5,-0.28) {1};
    \node[below] at (1.5,-0.95) {主视图};
  \end{scope}
  % 左视图
  \begin{scope}[shift={(5.2,0)}]
    \draw (0,0) rectangle (2,1);
    \draw (1,0) -- (1,1);
    \draw[<->] (0,-0.28) -- (1,-0.28); \node[below] at (0.5,-0.28) {1};
    \draw[<->] (1,-0.28) -- (2,-0.28); \node[below] at (1.5,-0.28) {1};
    \draw[<->] (2.35,0) -- (2.35,1); \node[right] at (2.35,0.5) {1};
    \node[below] at (1,-0.95) {左视图};
  \end{scope}
  % 俯视图
  \begin{scope}[shift={(0.2,-4.7)}]
    \draw (1,2) -- (2,2) -- (3,1) -- (2,0) -- (1,0) -- (0,1) -- cycle;
    \draw[<->] (1,2.3) -- (2,2.3); \node[above] at (1.5,2.3) {1};
    \draw[<->] (0.45,-0.3) -- (1,-0.3); \node[below] at (0.75,-0.3) {1};
    \draw[<->] (3.35,0) -- (3.35,1); \node[right] at (3.35,0.5) {1};
    \draw[<->] (3.35,1) -- (3.35,2); \node[right] at (3.35,1.5) {1};
    \node[below] at (1.5,-0.95) {俯视图};
  \end{scope}
\end{tikzpicture}
